% !TEX root = ../main.tex
% Numbers: global projection, household-budget-fix solution (6 Oct 2026), five nodes in the
% risk factor and last quarter's spread as a thirteenth state (125 collocation points); every
% economy solved to the collocation tolerance at every node. Headline welfare: Monte Carlo
% with every quarter cleared exactly (40 histories of 600 quarters, common shocks and default
% draws); the announcement component by exact clearing; the national-income decomposition by
% the solution's expectation operator. Results: results/WELFARE/<economy>/; drivers and the
% table generator in Claude files/global_projection/policy_rules/.
\section{Welfare and the design of the backstop}
\label{sec:welfare}

Along a history without default, a purchase of the $D$ bond hands the Eurosystem the excess return
the selling banks would have earned, and the state in which the insurance pays never arrives. We
therefore rank backstop rules by households' expected welfare over every history, including those in
which the backstop is activated and the sovereign defaults. We define the welfare criterion
and the rules, report welfare when the Eurosystem's profits are shared by the capital key, trace why
the vulnerable member loses, and then let that member keep the profits.

\subsection{The welfare criterion}
\label{sec:welfare-criterion}

\paragraph{Household welfare.}
Welfare is that of each country's representative household, with discount factor
$\hat\beta=1/(1+\bar r)$ and period utility $\log x_{X,t}$ ($\sigma_X=1$ in \eqref{eq:hh-objective}),
where $x_{X,t}=C_{X,t}-v(N_{X,t})$ is the GHH composite. As a function of the aggregate state $S_t$
and the default regime $d_t$,
\begin{equation}
V_X(S_t,d_t)=\log x_X(S_t,d_t)+\hat\beta\,\E_t\bigl[V_X(S_{t+1},d_{t+1})\bigr],
\qquad X\in\{D,F\},
\label{eq:welfare-bellman}
\end{equation}
where the expectation is the equilibrium's own, \eqref{eq:expectation}, with next quarter's regime
drawn at the priced probability $\pi^{d}(s_t)$. Default is not absorbing, so
\eqref{eq:welfare-bellman} averages over histories with any number of defaults.

\paragraph{Aggregation across countries.}
Union welfare weights every household equally,
\begin{equation}
W(S_t,d_t)=\frac{V_D(S_t,d_t)+\varkappa\,V_F(S_t,d_t)}{1+\varkappa}.
\label{eq:union-welfare}
\end{equation}
where $\varkappa$ is the size of $F$ relative to $D$. The countries are identical per capita in the steady state, so transfers between them move $W$ only
at second order: $W$ measures efficiency, $V_D$ and $V_F$ incidence.

\paragraph{Announcement state and consumption equivalents.}
Every rule is announced at $\mathcal S_0$, the stochastic rest point without a backstop (no
Eurosystem holdings, a 57\,bp spread, no default), so comparisons include the transition to the new
rest point. The consumption equivalent $\tau_X$ is the permanent transfer, as a share of rest-point
consumption $\bar C_X$, that would make households of $X$ without a backstop as well off as under
the rule,
\begin{equation}
\E_0\sum_{t=0}^{\infty}\hat\beta^{\,t}\log\bigl(x^{0}_{X,t}+\tau_X\bar C_X\bigr)
=V^{\text{rule}}_X(\mathcal S_0,0),
\label{eq:ce}
\end{equation}
where $x^{0}$ is the composite without a backstop. The union equivalent $\tau_U$, the same transfer in
both countries, is to first order the population average of $\tau_D$ and $\tau_F$.

\paragraph{Decompositions.}
We decompose welfare changes in two ways. The \emph{quiet history}, with no shocks or default,
isolates the \emph{announcement} because no purchases occur. The difference between total welfare
and the quiet history measures the contribution of \emph{crisis purchases}. We also decompose welfare
into output, terms of trade, investment, net exports and work through the national-income identity
$C_X=(Y_X-I_X-NX_X)/P^{c}_X$.

\subsection{Level and change rules}
\label{sec:welfare-rules}

The Eurosystem announces a target for its holdings of the $D$ sovereign, as a share of the
steady-state stock $\bar B_D$, that responds to the spread $\mathrm{sp}_t$ of \eqref{eq:spread-def}
(percentage points per year) and to its change:
\begin{equation}
\theta_t=\gamma_L\bigl(\mathrm{sp}_t-\bar s\bigr)+\gamma_\Delta\bigl(\mathrm{sp}_t-\mathrm{sp}_{t-1}\bigr).
\label{eq:gamma-rule}
\end{equation}
As under the spread cap of \cref{sec:policy}, the Eurosystem never sells and buys nothing in default:
\begin{equation}
m_t=\bar B_D\max\Bigl\{0,\ \min\{\theta_t,\bar\theta\}-(1-\delta_b)h_tM_{t-1}/\bar B_D\Bigr\},
\qquad m_t=0\ \text{if}\ d_t=1,
\label{eq:gamma-purchases}
\end{equation}
with book $M_t=(1-\delta_b)h_tM_{t-1}+m_t$. The reference spread $\bar s$ is the pre-crisis
rest-point spread, 57\,bp, and the issuer limit $\bar\theta=0.5$ is never approached. The coefficient
$\gamma_L$ is the share of the stock held per percentage point of spread above $\bar s$. With
$\gamma_L=0.15$, a crisis that takes the spread to 299\,bp without a backstop sets a target of 36\% of
the stock, close to the 37\% that the cap of \cref{sec:policy} buys in that state when
$\overline{\text{Spread}}=200$\,bp. We compare the level rule $R_L$ ($\gamma_\Delta=0$), the change
rule $R_\Delta$ ($\gamma_L=0$) and the combined rule $R_{L+\Delta}$. Under $R_L$ intervention persists
while financing conditions remain impaired; under $R_\Delta$ the Eurosystem arrests a disorderly rise
in the spread and stops once the market settles. The constant cap is the limit of $R_L$ as
$\gamma_L\to\infty$ with $\bar s$ equal to the cap.

\subsection{Welfare under the capital key}
\label{sec:welfare-results}

\Cref{tab:welfare} reports the welfare of each rule when the Eurosystem's profits and losses are
shared by the capital key ($\kappa_D=0.071$, $D$'s share of the profits; \cref{sec:policy}). Union welfare rises
under every rule, by $0.003$ to $0.011\%$ of consumption, and all of the gain comes from the
announcement. The promise lowers the rest-point lending spread, from 98 to 86\,bp under
$\gamma_L=0.1$, and so raises output before a single bond is bought. Crisis purchases are worth
slightly less than nothing to the union ($-0.0014$ to $-0.0001\%$): they insure $D$'s banks, but they
also set off the transfer problem of \cref{sec:welfare-channels}.

Households in $D$ lose $0.03$ to
$0.13\%$ of consumption and those in $F$ gain $0.007$ to $0.028\%$, both more under stronger rules.
Between a half and two-thirds of $D$'s loss comes from crisis purchases, because that is when the
Eurosystem collects the premia it then pays out mostly abroad. Under the level, change and combined
rules its expected profits are positive: it buys when the market spread exceeds the expected-loss
spread, and so collects the risk and liquidity premia the selling banks would otherwise have earned.
Expected discounted profits range from $0.5$ to $3.6\%$ of quarterly $D$ output, and $92.9\%$ of them
accrue to $F$ under the capital key. The 200\,bp cap buys less often but in larger amounts, which
leaves its expected profit too imprecisely estimated to distinguish from zero.

\begin{table}[tbp]
  \centering
  \begin{threeparttable}
  \caption{Welfare of announced backstops, profits and losses shared by capital key}
  \label{tab:welfare}
  \footnotesize\setlength{\tabcolsep}{3.5pt}
  \begin{tabular}{lcc rrr rr r rr}
    \toprule
    & & & \multicolumn{3}{c}{Welfare (\% of consumption)} & \multicolumn{2}{c}{Announcement}
      & & \multicolumn{2}{c}{Purchases} \\
    \cmidrule(lr){4-6}\cmidrule(lr){7-8}\cmidrule(lr){10-11}
    Rule & $\gamma_L$ & $\gamma_\Delta$ & Union & $D$ & $F$ & Union & $D$
      & Profit & Freq. & Book \\
    \midrule
    % generated by welfare_tables.py: capital key
    Level & 0.025 & 0 & 0.0031 & $-0.028$ & 0.0070 & 0.0032 & $-0.010$ & 0.9 & 10.7 & 1.2 \\
     & 0.05 & 0 & 0.0046 & $-0.047$ & 0.0110 & 0.0049 & $-0.016$ & 1.6 & 11.4 & 2.2 \\
     & 0.10 & 0 & 0.0067 & $-0.075$ & 0.0168 & 0.0073 & $-0.026$ & 2.5 & 11.3 & 3.9 \\
     & 0.15 & 0 & 0.0080 & $-0.097$ & 0.0211 & 0.0090 & $-0.033$ & 3.1 & 10.9 & 5.3 \\
     & 0.20 & 0 & 0.0091 & $-0.115$ & 0.0246 & 0.0105 & $-0.041$ & 3.6 & 10.5 & 6.5 \\
    Change & 0 & 0.025 & 0.0026 & $-0.029$ & 0.0066 & 0.0028 & $-0.012$ & 0.5 & 8.0 & 0.6 \\
     & 0 & 0.05 & 0.0053 & $-0.060$ & 0.0135 & 0.0056 & $-0.028$ & 1.0 & 9.7 & 1.3 \\
     & 0 & 0.10 & 0.0105 & $-0.127$ & 0.0276 & 0.0116 & $-0.067$ & 1.8 & 10.8 & 2.5 \\
    Combined & 0.05 & 0.025 & 0.0063 & $-0.068$ & 0.0156 & 0.0068 & $-0.027$ & 1.8 & 10.3 & 2.6 \\
     & 0.05 & 0.05 & 0.0082 & $-0.092$ & 0.0207 & 0.0090 & $-0.041$ & 2.0 & 9.8 & 3.0 \\
     & 0.10 & 0.025 & 0.0078 & $-0.088$ & 0.0199 & 0.0087 & $-0.034$ & 2.6 & 10.5 & 4.1 \\
     & 0.10 & 0.05 & 0.0095 & $-0.106$ & 0.0240 & 0.0106 & $-0.046$ & 2.6 & 9.9 & 4.4 \\
    Cap 200\,bp &  &  & 0.0095 & $-0.127$ & 0.0266 & 0.0102 & $-0.055$ & 1.3 & 1.5 & 4.6 \\
    \bottomrule
  \end{tabular}
  \begin{tablenotes}\footnotesize
  \item Welfare: consumption equivalents \eqref{eq:ce} at $\mathcal S_0$, by simulation with every
  quarter cleared exactly (40 histories of 600 quarters); union: population-weighted,
  \eqref{eq:union-welfare}. Announcement: the quiet history (no innovation, no default), cleared
  exactly. Profit: expected discounted profit of the Eurosystem over 600 quarters, \% of quarterly
  $D$ output. Purchases: share of quarters in which the Eurosystem buys more than 0.1\% of the stock,
  and its average book, \% of the stock. Sampling standard errors: $D$ between 0.0006 and 0.009, $F$ at
  most 0.001 (for the cap, 0.012 and 0.0014).
  \end{tablenotes}
  \end{threeparttable}
\end{table}

\subsection{Sources of the vulnerable member's loss}
\label{sec:welfare-channels}

Under the capital key $D$ suffers a transfer problem. For the level rule the present value of $D$'s
output barely moves (\cref{tab:welfare-channels}), but purchasing power moves abroad: the coupons on
the bonds the Eurosystem has acquired accrue mainly to $F$, and with home bias that shifts world
expenditure towards $F$ goods. With a trade elasticity of 0.5, the adjustment requires a large
deterioration in $D$'s terms of trade. For $\gamma_L=0.1$ it costs $0.11\%$ of consumption, partly
offset by $0.04\%$ less disutility of work. The transfer also makes $D$'s economy more volatile: under
$\gamma_L=0.1$ the standard deviation of output rises from $0.24\%$ without a backstop to $0.46\%$,
against $0.27$ to $0.29\%$ under the level and change rules when $D$ keeps the profits.

Keeping the profits in $D$ removes the transfer problem but not a second force, the fall in the
relative price of $D$'s good, which operates however the profits are allocated.
Relaxing the financial constraint expands $D$'s production: for $\gamma_L=0.1$ the extra output is
worth $0.48\%$ of consumption. Little of it stays in $D$. The lower relative price of $D$'s good
absorbs $0.24\%$ and the extra investment $0.23\%$, which leaves consumption barely higher, and the
additional hours cost $0.11\%$. At the new rest point output is $0.34\%$ higher, consumption slightly
lower and the price of the $F$ good $1.6\%$ higher. With a trade elasticity below one, much of the
gain from easier credit is exported to the core through the cheaper goods $D$ sells it.

\begin{table}[tbp]
  \centering
  \begin{threeparttable}
  \caption{Sources of the welfare change, \% of consumption}
  \label{tab:welfare-channels}
  \footnotesize
  \begin{tabular}{l rrrrr r}
    \toprule
    & Output & Terms of & Invest- & Net & Work & Total \\
    & & trade & ment & exports & & \\
    \midrule
    \multicolumn{7}{l}{\emph{$D$, profits shared by capital key}} \\
    Level, $\gamma_L=0.1$          & $-0.009$ & $-0.108$ & $-0.014$ & $-0.023$ & $0.043$ & $-0.111$ \\
    Change, $\gamma_\Delta=0.05$  & $-0.057$ & $-0.120$ & $0.006$ & $-0.029$ & $0.066$ & $-0.134$ \\
    Cap 200\,bp                    & $-0.096$ & $-0.202$ & $0.006$ & $-0.051$ & $0.112$ & $-0.230$ \\
    \multicolumn{7}{l}{\emph{$D$, profits kept by $D$ ($\kappa_D=1$)}} \\
    Level, $\gamma_L=0.1$          & $0.484$ & $-0.237$ & $-0.228$ & $0.008$ & $-0.110$ & $-0.084$ \\
    Cap 200\,bp                    & $0.534$ & $-0.262$ & $-0.252$ & $0.010$ & $-0.122$ & $-0.092$ \\
    \multicolumn{7}{l}{\emph{$F$, profits shared by capital key}} \\
    Level, $\gamma_L=0.1$          & $0.036$ & $0.013$ & $-0.013$ & $0.002$ & $-0.019$ & $0.019$ \\
    Cap 200\,bp                    & $0.065$ & $0.024$ & $-0.023$ & $0.004$ & $-0.035$ & $0.036$ \\
    \bottomrule
  \end{tabular}
  \begin{tablenotes}\footnotesize
  \item First-order decomposition of the change in the present value of the GHH composite,
  $x=C-v(N)$ with $C=(Y-I-NX)/P^{c}$, at $\mathcal S_0$, by the solution's expectation operator,
  expressed as a permanent consumption transfer. Output: the present value of output at
  rest-point prices; terms of trade: the change in its value from the price of the consumption
  basket; work: $-v(N)$. The totals for $D$ differ from the simulated ones of \cref{tab:welfare};
  the split between channels is the object of interest.
  \end{tablenotes}
  \end{threeparttable}
\end{table}

\subsection{Who bears the profits and losses}
\label{sec:welfare-key}

If the Eurosystem's profits and losses on $D$ bonds accrue to $D$'s treasury ($\kappa_D=1$), as the
November 2012 agreement to pass the Securities Markets Programme's income on Greek bonds back to
Greece partly arranged \citep{eurogroup2012}, the union gains two to three and a half times as much
from the same rule (\cref{tab:welfare-key}): $0.0187\%$ against $0.0067\%$ of consumption under
$\gamma_L=0.1$.

Almost all of the extra gain comes from a stronger announcement. Profits retained at home raise the
value banks place on future net worth, $\E[\Omega_D]$, and by \eqref{eq:mu-closed} a higher
$\E[\Omega_D]$ relaxes the constraint at given net worth and assets: a banker who values continuing
in business more has less to gain from diverting. Under $\gamma_L=0.1$, $\E[\Omega_D]$ at the rest
point rises from $1.159$ without a backstop to $1.173$, against $1.164$ under the capital key. The
multiplier falls from $0.0120$ to $0.0066$ ($0.0105$ under the capital key) and the lending spread to
54\,bp (86\,bp). Crisis purchases remain worth almost nothing to the union.

$D$ still loses under every rule, which is perhaps surprising once the transfer problem has gone.
Crisis purchases now insure it, worth up to $0.026\%$ of consumption, but the stronger announcement
expands its output further and passes more of the gain to $F$ through its relative price, at a cost
of $0.05$ to $0.19\%$. The insurance grows with the strength of the rule, so under the strongest level
rule $D$ loses $0.076\%$, against $0.115\%$ under the capital key.

\begin{table}[tbp]
  \centering
  \begin{threeparttable}
  \caption{Welfare when $D$ keeps the Eurosystem's profits and losses}
  \label{tab:welfare-key}
  \footnotesize\setlength{\tabcolsep}{3.5pt}
  \begin{tabular}{lcc rrr rr r rr}
    \toprule
    & & & \multicolumn{3}{c}{Welfare (\% of consumption)} & \multicolumn{2}{c}{Announcement}
      & & \multicolumn{2}{c}{Purchases} \\
    \cmidrule(lr){4-6}\cmidrule(lr){7-8}\cmidrule(lr){10-11}
    Rule & $\gamma_L$ & $\gamma_\Delta$ & Union & $D$ & $F$ & Union & $D$
      & Profit & Freq. & Book \\
    \midrule
    % generated by welfare_tables.py: national
    Level & 0.05 & 0 & 0.0158 & $-0.085$ & 0.0284 & 0.0150 & $-0.082$ & 1.7 & 11.0 & 2.3 \\
     & 0.10 & 0 & 0.0187 & $-0.098$ & 0.0332 & 0.0184 & $-0.104$ & 2.8 & 11.1 & 4.2 \\
     & 0.15 & 0 & 0.0186 & $-0.092$ & 0.0324 & 0.0188 & $-0.105$ & 3.8 & 11.0 & 5.9 \\
     & 0.20 & 0 & 0.0171 & $-0.076$ & 0.0287 & 0.0176 & $-0.094$ & 4.4 & 11.0 & 7.4 \\
    Change & 0 & 0.025 & 0.0092 & $-0.049$ & 0.0165 & 0.0091 & $-0.046$ & 0.5 & 8.0 & 0.7 \\
     & 0 & 0.05 & 0.0166 & $-0.098$ & 0.0309 & 0.0168 & $-0.101$ & 1.0 & 9.7 & 1.3 \\
     & 0 & 0.10 & 0.0264 & $-0.162$ & 0.0500 & 0.0271 & $-0.188$ & 1.8 & 10.7 & 2.5 \\
    Cap 200\,bp &  &  & 0.0214 & $-0.123$ & 0.0394 & 0.0211 & $-0.132$ & 2.4 & 1.8 & 6.0 \\
    \bottomrule
  \end{tabular}
  \begin{tablenotes}\footnotesize
  \item As \cref{tab:welfare}, with the Eurosystem's profits and losses on $D$ bonds remitted to $D$'s
  treasury ($\kappa_D=1$).
  \end{tablenotes}
  \end{threeparttable}
\end{table}

\subsection{Level against change}
\label{sec:welfare-level-change}

\Cref{fig:welfare-frontier} plots the union's gain against $D$'s loss for every rule. Under the
capital key (panel a) the level, change and combined rules and the 200\,bp cap lie on one frontier:
at equal loss for $D$, no rule departs from the level rules by more than $0.001$ points, within one
and a half standard errors. Along it the union gains $0.08$ to $0.11$ points of consumption per point
$D$ loses, less for stronger rules. Both the union's gain and $D$'s loss still rise at the strongest
rules we solve, so there is no interior optimum.

Rules at the same point of the frontier nonetheless behave differently. The change rule with
$\gamma_\Delta=0.05$ and the level rule with $\gamma_L=0.05$ cost $D$ similar amounts, but the change
rule holds a smaller book, $1.3\%$ of the stock against $2.2\%$, and destabilises $D$'s output less,
with a standard deviation of $0.29\%$ against $0.36\%$. It also compresses the tail of spreads less:
the 95th percentile is 189\,bp, against 170\,bp under the level rule and 193\,bp without a backstop.
The level rule keeps buying while spreads stay high, and so builds the larger book whose excess
return $D$ pays abroad.

When $D$ keeps the profits (panel b), the statistic the rule responds to matters for strong rules.
Weak level and change rules both yield $0.19$ union points per point of $D$'s loss. At a loss for $D$
of $0.098\%$, however, the level rule with $\gamma_L=0.1$ raises union welfare by $0.0187\%$ and the
change rule with $\gamma_\Delta=0.05$ by $0.0166\%$, four standard errors apart. Three quarters of the
difference is in the announcement, which lowers the rest-point lending spread to 54\,bp under the
level rule and 56\,bp under the change rule.

Beyond $\gamma_L=0.1$ the level rule's announcement stops growing: the rest-point lending spread
stays at 52 to 56\,bp and the union's gain at $0.017$ to $0.019\%$. Its crisis purchases, now
insurance for $D$, keep growing, so $D$'s loss falls. Perhaps surprisingly, the rule with
$\gamma_L=0.2$ costs $D$ less than that with $\gamma_L=0.05$ ($0.076$ against $0.085\%$) and gives $F$
as much. Its union gain of $0.017\%$ is one the change rule delivers only at a loss to $D$ of about
$0.10\%$.

Union gains above about $0.019\%$ require responding to the change in the spread, or the cap. With
$\gamma_\Delta=0.1$ the rest-point lending spread falls to 6\,bp, $D$'s bank constraint almost stops
binding at rest and the union gains $0.026\%$, but $D$ loses $0.162\%$: each further point of union
gain costs $D$ eight to nine points, about twice the average along the level rules.

\begin{figure}[tbp]
  \centering
  \includegraphics[width=\textwidth]{figures/fig_welfare_frontier}
  \caption{Union welfare gain against the loss of $D$'s households. Consumption equivalents in \% of
  consumption, by simulation with every quarter cleared exactly; bars are 95\% sampling intervals, for
  the union an upper bound.
  (a) Profits and losses shared by capital key. (b) Profits and losses kept by $D$. Labels give
  $\gamma_L$ or $\gamma_\Delta$; combined rules give both.}
  \label{fig:welfare-frontier}
\end{figure}

\subsection{Across households}
\label{sec:welfare-households}

The representative households hide how each member's loss or gain is shared within it. To measure
the incidence, the incomplete-markets households of \cref{sec:Households} re-solve their savings
problem under each economy's prices and aggregate law of motion, without feedback to market clearing
(\cref{tab:welfare-households}). Under $\gamma_L=0.1$ every quintile of $D$ loses about $0.10\%$ of
consumption ($0.08\%$ for the poorest). The loss is nearly flat because a terms-of-trade loss falls on
every household in proportion to the basket it buys; households above the poorest fifth lose slightly
more, consistent with a deposit rate about 1\,bp lower. The gain in $F$ rises with income, from
$0.010$ to $0.016\%$, consistent with a higher deposit rate. The change and combined rules distribute
their losses alike. When $D$ keeps the profits its deposit rate rises and its richest fifth loses
least ($0.05\%$ against about $0.07\%$). Because these households take the aggregate law of motion as
given, they misstate the level of $D$'s loss; only the shape of the incidence is informative.

\begin{table}[tbp]
  \centering
  \begin{threeparttable}
  \caption{Welfare across the household distribution, \% of consumption}
  \label{tab:welfare-households}
  \footnotesize\setlength{\tabcolsep}{4pt}
  \begin{tabular}{l rrrrr r r}
    \toprule
    & \multicolumn{6}{c}{$D$, by quintile of income} & $F$ \\
    \cmidrule(lr){2-7}
    Rule & Q1 & Q2 & Q3 & Q4 & Q5 & All & All \\
    \midrule
    % generated by welfare_tables.py: households
    Level, $\gamma_L=0.1$ & $-0.084$ & $-0.099$ & $-0.100$ & $-0.101$ & $-0.096$ & $-0.099$ & 0.0142 \\
    Change, $\gamma_\Delta=0.05$ & $-0.108$ & $-0.128$ & $-0.131$ & $-0.133$ & $-0.130$ & $-0.129$ & 0.0192 \\
    Combined, $(0.1,0.05)$ & $-0.152$ & $-0.179$ & $-0.183$ & $-0.185$ & $-0.178$ & $-0.180$ & 0.0257 \\
    Level, $\gamma_L=0.1$, $\kappa_D=1$ & $-0.072$ & $-0.078$ & $-0.070$ & $-0.066$ & $-0.047$ & $-0.070$ & 0.0274 \\
    \bottomrule
  \end{tabular}
  \begin{tablenotes}\footnotesize
  \item Incomplete-markets households of each country under the solved aggregate law of motion, at
  $\mathcal S_0$; Q1 is the poorest fifth by steady-state income (labour, transfers and interest).
  Uniform proportional rise in every household's GHH composite worth the welfare change, converted
  to consumption at the group's ratio of composite to consumption. Profits and losses shared by
  capital key unless stated.
  \end{tablenotes}
  \end{threeparttable}
\end{table}

\subsection{What the welfare analysis implies}
\label{sec:welfare-implications}

Under the capital key, the level, change and combined rules lie on essentially the same welfare
frontier. What matters is how much sovereign risk the Eurosystem ultimately absorbs rather than which
statistic of the spread triggers purchases. The allocation of profits has much larger welfare
consequences. Returning those profits to the vulnerable member substantially strengthens the
announcement effect and shifts the frontier outward, although the member continues to lose because
easier financing lowers the relative price of its own output. A Pareto-improving backstop therefore
requires the purchase rule to be accompanied by an explicit fiscal transfer.
